\documentclass[11pt,a4paper]{article}
\usepackage[T1]{fontenc}
\usepackage[utf8]{inputenc}
\usepackage{lmodern}
\usepackage{amsmath,amssymb}
\usepackage{booktabs}
\usepackage{microtype}
\usepackage{xcolor}
\usepackage[margin=2.5cm]{geometry}
\usepackage[colorlinks=true,linkcolor=blue!50!black,citecolor=blue!50!black,urlcolor=blue!50!black]{hyperref}

\newcommand{\Nc}{N_c}
\newcommand{\Ng}{N_g}
\newcommand{\MPl}{M_{\mathrm{Pl}}}
\newcommand{\sw}{\sin^2\theta_W}
\newcommand{\GeV}{\,\mathrm{GeV}}
\newcommand{\TeV}{\,\mathrm{TeV}}

\title{Three colours and three generations from perturbativity\\in 3-3-1 models}
\author{Sheldon Lindberg\\[2pt]\small Independent researcher\\[2pt]
\small Code and data: \url{https://github.com/sheldonlindberg-afk/STUR-Physics-Lab}}
\date{November 2026\\[4pt]\small Preprint, not peer reviewed}

\begin{document}
\maketitle

\begin{abstract}
In $SU(3)_C\times SU(3)_L\times U(1)_X$ (``3-3-1'') gauge models, anomaly cancellation
ties the number of fermion generations to the number of colours. We keep the number of
colours $\Nc$ general throughout, with every charge and anomaly coefficient computed as
a function of $\Nc$, and add a single requirement: no gauge coupling may become
non-perturbative below the Planck scale. The $SU(3)_L$ anomaly together with asymptotic
freedom of colour forces $\Nc$ to be odd and $\Ng=\Nc$, which makes Witten's global
anomaly condition automatic. $SU(3)_L$ running then excludes $\Nc\ge 7$ and $U(1)_X$
running excludes $\Nc=5$, so that $\Nc=\Ng=3$ is the only survivor, with quark charges
$\tfrac23$ and $-\tfrac13$. The selection holds at two loops in the gauge couplings with
two-loop Yukawa terms, with one-loop threshold matching across the heavy spectrum, for
3-3-1 breaking scales from 2 to $30\TeV$, and in both variants $\beta=\pm1/\sqrt3$. It
takes the measured electroweak couplings at the TeV scale as input and does not predict
them; in the $\beta=+1/\sqrt3$ variant the $\Nc=3$ margin is a factor of about 2 in the
pole scale. Requiring the scalar sector itself to remain perturbative and bounded below
up to $\MPl$ then bounds the exotic-quark masses relative to the bilepton mass. All
numbers are reproduced by open code with regression tests.
\end{abstract}

\section{Introduction}

In 3-3-1 models~\cite{Pisano:1992bxx,Frampton:1992wt} the $SU(3)_L$ anomaly cancels
only between generations, which links the number of generations to the number of
colours: $\Ng$ must be a multiple of $\Nc$~\cite{Pisano:1996ht,Doff:2000pc,Doff:2000tx}.
With the number of light neutrinos taken from experiment, Frampton noted that this
singles out $\Nc=3$~\cite{Frampton:1994rt}. Here no generation or neutrino number is
taken from experiment. We ask instead which $\Nc$ is compatible with all gauge couplings
remaining perturbative up to the Planck scale, given the measured electroweak couplings
at the TeV scale. The answer is $\Nc=\Ng=3$ only.

Section~\ref{sec:setup} sets up the model for general $\Nc$. Section~\ref{sec:chain}
gives the four-step selection; Sections~\ref{sec:robust}--\ref{sec:scalar} test its
robustness, including one-loop threshold matching and the scalar sector. The
conditions under which it holds are collected in Section~\ref{sec:caveats}, and
Section~\ref{sec:prior} compares with prior work.

\section{Setup for general $N_c$}\label{sec:setup}

The electric charge is $Q=T_3+\beta T_8+X$. The two variants without low-scale Landau
poles are $\beta=\pm1/\sqrt3$; the minimal $\beta=-\sqrt3$ model is excluded by atomic
parity violation and sits at the LHC $Z'$ bound (Section~\ref{sec:tests}). For colour group $SU(\Nc)$ we take $\Ng$ generations.
Each has left-handed lepton triplets and quark triplets or anti-triplets, the third member
of each quark multiplet being an exotic quark; $n$ of the $\Ng$ quark generations sit in
the anti-triplet.

Cancelling the perturbative anomalies, with the Yukawa couplings in place, gives
\begin{equation}
  Q_u=\frac{\Nc+1}{2\Nc},\qquad Q_d=\frac{1-\Nc}{2\Nc},
\end{equation}
which reduce to $\tfrac23$ and $-\tfrac13$ at $\Nc=3$. Every gauge anomaly ($X^3$,
grav$^2X$, $SU(3)_L^2X$, $SU(\Nc)^2X$, $SU(3)_L^3$) then vanishes identically in $\Nc$,
both in the Standard Model and in the full 3-3-1 content; this was checked in exact
rational arithmetic for $\Nc=3,5,7$. Perturbative anomalies therefore fix the charges
given $\Nc$, but not $\Nc$ itself.

The $U(1)_X$ charges for $\beta=-1/\sqrt3$ are $X=Q_u-\tfrac13$ (quark triplet),
$Q_d+\tfrac13$ (anti-triplet), $-\tfrac13$ (leptons), and $-\tfrac13,+\tfrac23,-\tfrac13$
for the scalars $\eta,\rho,\chi$, fixed by the Yukawas and independent of $\Nc$. For
$\beta=+1/\sqrt3$: leptons $-\tfrac23$, quark triplet $Q_u-\tfrac23$, anti-triplet
$Q_d+\tfrac23$, scalars $\tfrac13,\tfrac13,-\tfrac23$.

We require a non-abelian colour group, $\Nc\ge2$: for $\Nc=1$ there is no colour gauge
field at all, and the asymptotic-freedom condition used below is vacuous.

\section{The selection}\label{sec:chain}

\paragraph{Step 1: $\Ng=\Nc$.}
$SU(3)_L$ anomaly freedom requires equal numbers of triplets and anti-triplets,
\begin{equation}\label{eq:anomaly}
  \Ng(1+\Nc)=2n\Nc .
\end{equation}
For odd $\Nc$, $\gcd(\Nc,\Nc+1)=1$, so $\Nc\mid\Ng$ and $\Ng=k\Nc$. With three quark
flavours per generation, asymptotic freedom of $SU(\Nc)$ requires
$11\Nc/3-2\Ng>0$, i.e.\ $k<11/6$, so $k=1$. For even $\Nc$, $\gcd(\Nc+1,2\Nc)=1$, so
$2\Nc\mid\Ng$; but $2\Nc>11\Nc/6$ violates asymptotic freedom. Hence $\Nc$ is odd and
$\Ng=\Nc$, with no further input. At $\Nc=3$ this reproduces the standard assignment of
two quark generations in the anti-triplet and one in the triplet.

\paragraph{Step 2: Witten's anomaly is implied.}
In the Standard Model, Witten's global $SU(2)$ anomaly~\cite{Witten:1982fp} requires an
even number of doublets, $\Nc+1$ per generation. In 3-3-1, $SU(2)_L\subset SU(3)_L$ and
$\pi_4(SU(3))=0$, so the UV group imposes no global anomaly of its own. Step~1 already
forces $\Nc$ odd, and the doublet count $\Ng(1+\Nc)=2n\Nc$ is then even automatically.
The spin--charge relation $(-1)^F=\exp(i\pi\Nc(B-L))$ gives the same condition as a
consequence rather than an input (all 29 solutions of~\eqref{eq:anomaly} up to $\Nc=60$
checked).

\paragraph{Step 3: $SU(3)_L$ running excludes $N_c\ge7$.}
With $\Ng=\Nc$ there are $\Nc(1+\Nc)$ Weyl triplets. The one-loop coefficient
$b_L=11-\Nc(1+\Nc)/3-(\text{scalars})$, run from $1/\alpha_L(4\TeV)=31.5$, gives
Table~\ref{tab:su3L}. Adding the minimal model's scalar sextet does not change the
verdict.

\begin{table}[h]\centering
\begin{tabular}{ccl}
\toprule
$\Nc$ & $b_L$ (3 scalar triplets) & Landau pole \\
\midrule
3 & $+6.5$ & none (asymptotically free)\\
5 & $+0.5$ & none\\
7 & $-8.2$ & $1.3\times10^{14}\GeV$, below $\MPl$\\
9 & $-19.5$ & $10^{8}\GeV$\\
\bottomrule
\end{tabular}
\caption{One-loop $SU(3)_L$ running with $\Ng=\Nc$.}\label{tab:su3L}
\end{table}

\paragraph{Step 4: $U(1)_X$ running excludes $N_c=5$.}
Two-loop gauge running from the 3-3-1 scale starts at $1/\alpha_X(4\TeV)=83.76$; we
define ``blow-up'' as $g=4\pi$. The results are in Table~\ref{tab:x}; the Planck scale
is $1.2\times10^{19}\GeV$. For $\beta=-1/\sqrt3$ a pole exactly at $\MPl$ needs
$a_X=14.76$, so the one-loop margins are $+6.09$ for $\Nc=3$ and $-4.57$ for $\Nc=5$.
Two-loop gauge terms move the blow-up scale by a factor of about 3.

\begin{table}[h]\centering
\begin{tabular}{cccc}
\toprule
$\Nc=\Ng$ & one-loop $a_X$ ($\beta=-1/\sqrt3$;\ $+1/\sqrt3$) & blow-up, $\beta=-1/\sqrt3$ & blow-up, $\beta=+1/\sqrt3$\\
\midrule
3 & $26/3=8.67$;\ $14.00$ & $\mathbf{3.0\times10^{29}\GeV}$ & $\mathbf{4.0\times10^{19}\GeV}$\\
5 & $19.33$;\ $27.33$ & $1.3\times10^{15}\GeV$ & $5.5\times10^{11}\GeV$\\
7 & $34.44$;\ $45.11$ & $9.7\times10^{9}\GeV$ & $3.1\times10^{8}\GeV$\\
\bottomrule
\end{tabular}
\caption{Two-loop $U(1)_X$ blow-up scales.}\label{tab:x}
\end{table}

\paragraph{Two-loop Yukawa terms.}
A Yukawa coupling $y\,\phi\psi_1\psi_2$ contributes
$-g^3y^2\,n_{\text{terms}}\,[C_2(\psi_1)+C_2(\psi_2)]/d(G)/(16\pi^2)^2$ to each gauge
$\beta$-function~\cite{Machacek:1983tz,Luo:2002ti}; this reproduces the Standard Model
top-quark coefficients $(17/10,\,3/2,\,2)$ exactly. We include the top Yukawa
($y_t=0.83$ at $4\TeV$) and one Yukawa $y_J$ for every exotic quark, and for every heavy
lepton when $\beta=+1/\sqrt3$, held constant up to the blow-up. Real Yukawas decrease
with energy, so this overstates their effect and makes the $\Nc=5$ exclusion
conservative. Yukawas raise every pole by at most a factor of about 5
(Table~\ref{tab:yuk}), and no constant $y_J\le6$ lifts $\Nc=5$ to the Planck scale.

\begin{table}[h]\centering
\begin{tabular}{ccc}
\toprule
$y_J$ & $\beta=-1/\sqrt3$: $\Nc=3$ / $\Nc=5$ & $\beta=+1/\sqrt3$: $\Nc=3$ / $\Nc=5$\\
\midrule
no Yukawas & $3.0\times10^{29}$ / $1.3\times10^{15}\GeV$ & $4.0\times10^{19}$ / $5.5\times10^{11}\GeV$\\
0 (top only) & $3.4\times10^{29}$ / $1.3\times10^{15}$ & $4.2\times10^{19}$ / $5.6\times10^{11}$\\
1 & $4.7\times10^{29}$ / $1.6\times10^{15}$ & $6.2\times10^{19}$ / $6.8\times10^{11}$\\
2 & $1.2\times10^{30}$ / $2.8\times10^{15}$ & $2.2\times10^{20}$ / $1.3\times10^{12}$\\
\bottomrule
\end{tabular}
\caption{Blow-up scales with two-loop Yukawa contributions.}\label{tab:yuk}
\end{table}

\section{Robustness}\label{sec:robust}

The two-loop machinery reproduces the Standard Model one-loop coefficients
$(41/10,-19/6,-7)$ and all nine two-loop coefficients exactly. Table~\ref{tab:robust}
varies the 3-3-1 scale and the perturbativity criterion.

\begin{table}[h]\centering\small
\begin{tabular}{lll}
\toprule
Check & $\beta=-1/\sqrt3$: $\Nc=3$ vs $5$ & $\beta=+1/\sqrt3$: $\Nc=3$ vs $5$\\
\midrule
3-3-1 scale $2$--$30\TeV$ & $(2.7$--$3.7)\times10^{29}$ vs $(0.8$--$4.5)\times10^{15}$ & $(2.9$--$9.9)\times10^{19}$ vs $(3.3$--$24)\times10^{11}$\\
criterion $\alpha=1$ & $1.8\times10^{29}$ vs $1.0\times10^{15}$ & $2.9\times10^{19}$ vs $4.7\times10^{11}$\\
criterion $g=1$ (strict) & $5.3\times10^{25}$ vs $2.7\times10^{13}$ & $2.0\times10^{17}$ vs $3.6\times10^{10}$ ($\Nc=3$ fails)\\
\bottomrule
\end{tabular}
\caption{Blow-up scales in GeV under variations.}\label{tab:robust}
\end{table}

For $\beta=-1/\sqrt3$, $\Nc=\Ng=3$ is selected under every criterion, with its blow-up
6--10 orders of magnitude above $\MPl$ while $\Nc=5$ fails 4--6 orders below. For
$\beta=+1/\sqrt3$, 3 is selected under the standard criteria, but its pole lies only a
factor 2--3 above $\MPl$ without Yukawas (3--5 with them); under the strict $g=1$
criterion, $\Nc=3$ fails as well.

\section{One-loop threshold matching}\label{sec:threshold}

The breaking $SU(3)_L\times U(1)_X\to SU(2)_L\times U(1)_Y$ is matched at one loop in
$\overline{\text{MS}}$~\cite{Weinberg:1980wa,Hall:1980kf}:
\begin{equation}
\frac{1}{\alpha_{\text{low}}(\mu)}=\frac{1}{\alpha_{\text{high}}(\mu)}-\frac{\lambda}{12\pi},\quad
\lambda=\sum_V t^2\Bigl(1-21\ln\frac{M_V}{\mu}\Bigr)+8\sum_{\text{Dirac}}t^2\ln\frac{M_F}{\mu}
+\sum_{\text{real scalars}}t^2\ln\frac{M_S}{\mu},
\end{equation}
with $1/\alpha_X+\beta^2/\alpha_L$ as the high-energy side for $U(1)_Y$. The heavy
spectrum is: one complex $SU(2)$ doublet of bilepton vectors ($|Y|=\tfrac12$, mass
$M_V$); the $Z'$ ($M_{Z'}=1.225\,M_V$); $\Nc$ exotic quarks (plus $\Nc$ heavy leptons for
$\beta=+1/\sqrt3$) of mass $r_JM_V$, whose Yukawa $y_J=r_Jg_L/\sqrt2$ enters at two loops;
and nine heavy real scalars of mass $r_SM_V$ (a second doublet, a charged singlet and three
neutral fields), with the option of a light second doublet.

As a check, at one loop the $\Nc=3$ pole must not depend on the matching scale once the
spectrum is fixed: moving $\mu$ from $M_V/3$ to $3M_V$ changes $\ln(\text{pole})$ by less
than $10^{-9}$, against a factor 2.6 without thresholds. $\Nc=5$ cannot pass this check,
because the low-energy theory is the three-colour Standard Model; its pole moves by a
factor 1.5 with $\mu$, and we take its most favourable $\mu$.

Scanning $r_J=0.5$--$3$, $r_S=\tfrac13$--$3$, $M_V=2$--$10\TeV$ and a heavy or light
second doublet (two-loop running, Yukawas included) gives Table~\ref{tab:thr}. The worst
case for $\Nc=3$ (lightest exotics, light scalars, light second doublet, $M_V=2\TeV$) puts
the $Z'$ at $2.45\TeV$, already below the conservative LHC limits of
Section~\ref{sec:tests}. At two loops, moving $\mu$ from $M_V/2$
to $2M_V$ changes the pole by 1\%.

\begin{table}[h]\centering
\begin{tabular}{lccc}
\toprule
& $\Nc=3$, worst case & $\Nc=5$, best case & verdict\\
\midrule
$\beta=-1/\sqrt3$ & $2.4\times10^{29}\GeV$ ($2\times10^{10}\MPl$) & $7.2\times10^{15}\GeV$ & only $\Nc=3$\\
$\beta=+1/\sqrt3$ & $2.0\times10^{19}\GeV$ ($1.64\,\MPl$) & $3.8\times10^{12}\GeV$ & only $\Nc=3$\\
$\beta=+1/\sqrt3$, $M_V\ge4\TeV$ & $2.23\,\MPl$ & --- & only $\Nc=3$\\
\bottomrule
\end{tabular}
\caption{Selection with one-loop threshold matching over the scanned spectrum.}\label{tab:thr}
\end{table}

The selection holds for any cutoff $\Lambda$ between the best $\Nc=5$ pole and the worst
$\Nc=3$ pole. Over all criteria ($g=4\pi$ or $\alpha=1$, with or without Yukawas), any
$\Lambda$ from $7.2\times10^{15}$ to $1.4\times10^{19}\GeV$ works. This window contains
$\MPl$, the reduced Planck mass and the conventional unification scale.

\section{The scalar sector to the Planck scale}\label{sec:scalar}

The selection requires every coupling of 3-3-1 at $\Nc=3$ to stay perturbative to $\MPl$,
so the scalar sector must pass the same test. We implement the general one-loop quartic
and Yukawa RGEs~\cite{Machacek:1983fi,Machacek:1984zw,Luo:2002ti} on the 18 real scalars;
the engine reproduces the Standard Model $\beta_\lambda$ and $\beta_{y_t}$ exactly. The
nine quartics allowed by the model's discrete symmetry close under the RGEs (residual
$10^{-15}$), and the potential is gauge invariant to $10^{-16}$. The light-Higgs quartic is
matched at $4\TeV$ to $\lambda_{\text{SM}}=0.076$ at tree level in the decoupling limit.
Survival means every $|\lambda|<4\pi$, every Yukawa $<\sqrt{4\pi}$, and a potential
bounded below, checked numerically at eight scales up to $\MPl$.

3-3-1 at $\Nc=3$ survives in both variants, in a corner of parameter space: 17 of 250
viable points ($\beta=-1/\sqrt3$) and 13 of 250 ($\beta=+1/\sqrt3$) under small priors,
none with $O(1)$ quartics. In the survivors the extra scalars cure the Standard Model's
metastability. Consistency to $\MPl$ then requires the exotic Yukawa at $4\TeV$ to be
below 0.80 ($\beta=-1/\sqrt3$) or 0.62 ($\beta=+1/\sqrt3$), i.e.\ exotic-quark masses below
$1.79\,M_V$ or $1.40\,M_V$; the top's doublet must carry $c^2\gtrsim0.68$ of the
electroweak vev squared; and all quartics at $4\TeV$ must be $\lesssim0.37$. These
statements are one-loop, with tree-level matching, and the search is a random sampling
rather than a proof that the surviving region is unique.

\section{Conditions and caveats}\label{sec:caveats}

\begin{enumerate}
\item \emph{The framework.} Everything assumes 3-3-1 with $\beta=\pm1/\sqrt3$.
\item \emph{Perturbativity to the Planck scale} is a standard criterion, not a theorem. A
pole below $\MPl$ signals new physics rather than an inconsistency.
\item \emph{Measured couplings are input.} Every $\Nc$ is run from $\sw=0.2504$,
$1/\alpha_{\text{em}}=125.7$ and $\alpha_s$ at $4\TeV$; in 3-3-1,
$1/\alpha_X=(1-4s^2/3)/\alpha_{\text{em}}$. For $\beta=-1/\sqrt3$, $\Nc=3$ reaches $\MPl$
for $\sw<0.449$, while $\Nc=5$ does so only for $0.066<\sw<0.074$, a factor 3.1--3.4
below the measured value. For $\beta=+1/\sqrt3$, $\Nc=5$ fails at every weak angle, and
$\Nc=3$ survives for $\sw<0.266$ (tree-level matching; $+0.010$ to $+0.035$ in margin with
thresholds). The precise statement is therefore: \emph{given the measured electroweak
couplings}, only $\Nc=3$ stays perturbative to $\MPl$.
\item \emph{Minimality is not explained.} Steps 1--2 leave $\{3,5,7,\dots\}$; the running
does the selecting.
\item \emph{No numbers follow.} The structure adds about eleven scalar parameters beyond
the Standard Model's two and blocks gauge unification ($SU(3)_C$ and $SU(3)_L$ never meet).
\end{enumerate}

\section{Relation to prior work}\label{sec:prior}

That $\Ng$ is a multiple of $\Nc$ in 3-3-1 is known~\cite{Pisano:1992bxx,Frampton:1992wt,Pisano:1996ht,Doff:2000pc}.
Frampton obtained $\Nc=3$ by inputting three light neutrinos~\cite{Frampton:1994rt}.
Charge quantization for arbitrary $\Nc$ was discussed in~\cite{Geng:2000pf}, generation
counting for arbitrary $\beta$ in~\cite{Diaz:2004fs}, and the $\Nc$-dependence of the
charges in~\cite{Bar:2001qk}. The closest predecessor is Doff and Pisano~\cite{Doff:2000tx},
who keep $n_c$ general and find $\Ng=k\,n_c$ for any $k\ge1$; they do not use asymptotic
freedom to set $k=1$, nor running to choose $n_c$. Landau poles in 3-3-1 at $\Nc=3$ are
studied in~\cite{Dias:2004dc,Barela:2023oyp,Morisi:2025svd}.

The step taken here is to keep $\Nc$ general throughout (charges, anomalies, scalars,
two-loop running, thresholds, scalar sector) without taking the number of generations or
neutrinos from experiment. Searches of arXiv and INSPIRE, including every paper citing
the founding 3-3-1 papers, Frampton~\cite{Frampton:1994rt} and Doff--Pisano~\cite{Doff:2000tx},
found no prior use of asymptotic freedom plus running to fix $\Nc$. This is not proof of
novelty, and comments from experts are welcome.

\section{Experimental tests}\label{sec:tests}

The $\beta=\pm1/\sqrt3$ models predict a $Z'$, singly or doubly charged bilepton gauge
bosons and exotic quarks, which the LHC and future colliders will find or exclude.
Current dilepton limits on the $Z'$ depend on the variant and on its decays. For
$\beta=-1/\sqrt3$, a recast of $13\TeV$ ATLAS data (36.1\,fb$^{-1}$) gives
$M_{Z'}>4\TeV$ if the $Z'$ decays only to Standard Model particles~\cite{Lindner:2016bgg};
with decays into exotic quarks and heavy leptons open, the bound weakens to
$2.64\TeV$~\cite{deJesus:2020ngn}. For $\beta=+1/\sqrt3$ the most recent dedicated recast we
found uses $8\TeV$ data and gives $M_{Z'}\gtrsim2.9\TeV$~\cite{Salazar:2015gxa}; it is
likely superseded by $13\TeV$ data, which we have not recast.

Consistency to $\MPl$ adds the relations of Section~\ref{sec:scalar}: exotic quarks
lighter than $1.79\,M_V$ or $1.40\,M_V$, $M_{Z'}=1.225\,M_V$, and the top's doublet
carrying $\gtrsim70\%$ of the electroweak vev squared. A discovery violating these would
break the selection without excluding 3-3-1.

In the minimal $\beta=-\sqrt3$ model, $\sw$ reaches $\tfrac14$ at $3.68\TeV$ at one loop,
where $g_X$ has its Landau pole, so its $Z'$ must lie below about $3.7\TeV$. The LHC
bound in this model is $M_{Z'}>3.7\TeV$ ($13\TeV$, 36.1\,fb$^{-1}$, exotic decays
included)~\cite{Nepomuceno:2019eaz}, and atomic parity violation in caesium requires
$M_{Z'}\gtrsim4.2$--$5.4\TeV$~\cite{Long:2018fud}. The minimal model is therefore excluded
unless higher-order effects raise its pole by more than about 14\%.

Finally, $\Nc$ can be checked directly at low energy. With the $\Nc$-dependent charges of
Section~\ref{sec:setup} and the $O(\alpha_s)$ correction, $R=\sigma(e^+e^-\to
\text{hadrons})/\sigma(e^+e^-\to\mu^+\mu^-)$ below charm is
$\Nc(Q_u^2+2Q_d^2)(1+3C_F\alpha_s/4\pi)$, which is $2.173$ at $\Nc=3$. KEDR's weighted
average over 1.84--3.05\,GeV is $R=2.225\pm0.020\pm0.047$~\cite{KEDR:2016hjj}, so $\Nc=3$
agrees to $1.0\sigma$, while $\Nc=2$ and $5$ miss by $15\sigma$ and $33\sigma$. This check
is independent of the running argument; it is of course not new, but it confirms that the
$\Nc$-dependent charges reproduce the measured value.

\section*{Reproducibility}

Each number above is produced by a script in \texttt{src/m331\_*.py}, with reference
output in \texttt{results/m331/} and regression tests in \texttt{tests/test\_331\_*.py},
run on every commit.

\begin{thebibliography}{99}
\bibitem{Pisano:1992bxx} F.~Pisano and V.~Pleitez, \emph{An SU(3)$\times$U(1) model for electroweak interactions}, Phys.\ Rev.\ D \textbf{46}, 410 (1992), \href{https://doi.org/10.1103/PhysRevD.46.410}{doi:10.1103/PhysRevD.46.410}, \href{https://arxiv.org/abs/hep-ph/9206242}{hep-ph/9206242}.
\bibitem{Frampton:1992wt} P.~H.~Frampton, \emph{Chiral dilepton model and the flavor question}, Phys.\ Rev.\ Lett.\ \textbf{69}, 2889 (1992), \href{https://doi.org/10.1103/PhysRevLett.69.2889}{doi:10.1103/PhysRevLett.69.2889}.
\bibitem{Pisano:1996ht} F.~Pisano, \emph{A simple solution for the flavor question}, Mod.\ Phys.\ Lett.\ A \textbf{11}, 2639 (1996), \href{https://doi.org/10.1142/S0217732396002630}{doi:10.1142/S0217732396002630}, \href{https://arxiv.org/abs/hep-ph/9609358}{hep-ph/9609358}.
\bibitem{Doff:2000pc} A.~Doff and F.~Pisano, \emph{Minimal extended flavor groups, matter fields chiral representations, and the flavor question}, Mod.\ Phys.\ Lett.\ A \textbf{15}, 1471 (2000), \href{https://doi.org/10.1142/S021773230000178X}{doi:10.1142/S021773230000178X}, \href{https://arxiv.org/abs/hep-ph/0007306}{hep-ph/0007306}.
\bibitem{Doff:2000tx} A.~Doff and F.~Pisano, \emph{On the color simple group from chiral electroweak gauge groups}, \href{https://arxiv.org/abs/hep-ph/0011087}{hep-ph/0011087} (2000).
\bibitem{Frampton:1994rt} P.~H.~Frampton, \emph{The third family is different}, in Proc.\ PASCOS 94, p.~63, \href{https://arxiv.org/abs/hep-ph/9409331}{hep-ph/9409331}.
\bibitem{Witten:1982fp} E.~Witten, \emph{An SU(2) anomaly}, Phys.\ Lett.\ B \textbf{117}, 324 (1982), \href{https://doi.org/10.1016/0370-2693(82)90728-6}{doi:10.1016/0370-2693(82)90728-6}.
\bibitem{Machacek:1983tz} M.~E.~Machacek and M.~T.~Vaughn, \emph{Two-loop renormalization group equations in a general quantum field theory. 1. Wave function renormalization}, Nucl.\ Phys.\ B \textbf{222}, 83 (1983), \href{https://doi.org/10.1016/0550-3213(83)90610-7}{doi:10.1016/0550-3213(83)90610-7}.
\bibitem{Machacek:1983fi} M.~E.~Machacek and M.~T.~Vaughn, \emph{\dots\ 2. Yukawa couplings}, Nucl.\ Phys.\ B \textbf{236}, 221 (1984), \href{https://doi.org/10.1016/0550-3213(84)90533-9}{doi:10.1016/0550-3213(84)90533-9}.
\bibitem{Machacek:1984zw} M.~E.~Machacek and M.~T.~Vaughn, \emph{\dots\ 3. Scalar quartic couplings}, Nucl.\ Phys.\ B \textbf{249}, 70 (1985), \href{https://doi.org/10.1016/0550-3213(85)90040-9}{doi:10.1016/0550-3213(85)90040-9}.
\bibitem{Luo:2002ti} M.-x.~Luo, H.-w.~Wang and Y.~Xiao, \emph{Two-loop renormalization group equations in general gauge field theories}, Phys.\ Rev.\ D \textbf{67}, 065019 (2003), \href{https://doi.org/10.1103/PhysRevD.67.065019}{doi:10.1103/PhysRevD.67.065019}, \href{https://arxiv.org/abs/hep-ph/0211440}{hep-ph/0211440}.
\bibitem{Weinberg:1980wa} S.~Weinberg, \emph{Effective gauge theories}, Phys.\ Lett.\ B \textbf{91}, 51 (1980), \href{https://doi.org/10.1016/0370-2693(80)90660-7}{doi:10.1016/0370-2693(80)90660-7}.
\bibitem{Hall:1980kf} L.~J.~Hall, \emph{Grand unification of effective gauge theories}, Nucl.\ Phys.\ B \textbf{178}, 75 (1981), \href{https://doi.org/10.1016/0550-3213(81)90498-3}{doi:10.1016/0550-3213(81)90498-3}.
\bibitem{Geng:2000pf} C.~Q.~Geng, \emph{Anomaly, charge quantization and family}, in Proc.\ Int.\ Workshop on Nonperturbative Methods and Lattice QCD (2000), p.~294, \href{https://doi.org/10.1142/9789812811370_0030}{doi:10.1142/9789812811370\_0030}, \href{https://arxiv.org/abs/hep-ph/0101329}{hep-ph/0101329}.
\bibitem{Diaz:2004fs} R.~A.~Diaz, R.~Martinez and F.~Ochoa, \emph{SU(3)$_c\times$SU(3)$_L\times$U(1)$_X$ models for $\beta$ arbitrary and families with mirror fermions}, Phys.\ Rev.\ D \textbf{72}, 035018 (2005), \href{https://doi.org/10.1103/PhysRevD.72.035018}{doi:10.1103/PhysRevD.72.035018}, \href{https://arxiv.org/abs/hep-ph/0411263}{hep-ph/0411263}.
\bibitem{Bar:2001qk} O.~B\"ar and U.-J.~Wiese, \emph{Can one see the number of colors?}, Nucl.\ Phys.\ B \textbf{609}, 225 (2001), \href{https://doi.org/10.1016/S0550-3213(01)00288-7}{doi:10.1016/S0550-3213(01)00288-7}, \href{https://arxiv.org/abs/hep-ph/0105258}{hep-ph/0105258}.
\bibitem{Dias:2004dc} A.~G.~Dias, R.~Martinez and V.~Pleitez, \emph{Concerning the Landau pole in 3-3-1 models}, Eur.\ Phys.\ J.\ C \textbf{39}, 101 (2005), \href{https://doi.org/10.1140/epjc/s2004-02083-0}{doi:10.1140/epjc/s2004-02083-0}, \href{https://arxiv.org/abs/hep-ph/0407141}{hep-ph/0407141}.
\bibitem{Barela:2023oyp} M.~W.~Barela, \emph{On the 3-3-1 Landau pole}, Nucl.\ Phys.\ B \textbf{1000}, 116475 (2024), \href{https://doi.org/10.1016/j.nuclphysb.2024.116475}{doi:10.1016/j.nuclphysb.2024.116475}, \href{https://arxiv.org/abs/2305.05066}{arXiv:2305.05066}.
\bibitem{Morisi:2025svd} S.~Morisi, G.~Ricciardi and G.~P.~Perdon\`a, \emph{Rescuing 331 bileptons from the Landau pole}, Phys.\ Rev.\ D \textbf{112}, 035041 (2025), \href{https://doi.org/10.1103/c7jb-lbsx}{doi:10.1103/c7jb-lbsx}, \href{https://arxiv.org/abs/2505.15785}{arXiv:2505.15785}.
\bibitem{Lindner:2016bgg} M.~Lindner, M.~Platscher and F.~S.~Queiroz, \emph{A call for new physics: the muon anomalous magnetic moment and lepton flavor violation}, Phys.\ Rept.\ \textbf{731}, 1 (2018), \href{https://doi.org/10.1016/j.physrep.2017.12.001}{doi:10.1016/j.physrep.2017.12.001}, \href{https://arxiv.org/abs/1610.06587}{arXiv:1610.06587}.
\bibitem{deJesus:2020ngn} \'A.~S.~de Jesus, S.~Kovalenko, F.~S.~Queiroz, C.~A.~de S.~Pires and Y.~S.~Villamizar, \emph{Dead or alive? Implications of the muon anomalous magnetic moment for 3-3-1 models}, Phys.\ Lett.\ B \textbf{809}, 135689 (2020), \href{https://doi.org/10.1016/j.physletb.2020.135689}{doi:10.1016/j.physletb.2020.135689}, \href{https://arxiv.org/abs/2003.06440}{arXiv:2003.06440}.
\bibitem{Salazar:2015gxa} C.~Salazar, R.~H.~Benavides, W.~A.~Ponce and E.~Rojas, \emph{LHC constraints on 3-3-1 models}, JHEP \textbf{07}, 096 (2015), \href{https://doi.org/10.1007/JHEP07(2015)096}{doi:10.1007/JHEP07(2015)096}, \href{https://arxiv.org/abs/1503.03519}{arXiv:1503.03519}.
\bibitem{Nepomuceno:2019eaz} A.~Nepomuceno and B.~Meirose, \emph{Limits on 331 vector bosons from LHC proton collision data}, Phys.\ Rev.\ D \textbf{101}, 035017 (2020), \href{https://doi.org/10.1103/PhysRevD.101.035017}{doi:10.1103/PhysRevD.101.035017}, \href{https://arxiv.org/abs/1911.12783}{arXiv:1911.12783}.
\bibitem{Long:2018fud} H.~N.~Long, N.~V.~Hop, L.~T.~Hue and N.~T.~T.~Van, \emph{Constraining heavy neutral gauge boson $Z'$ in the 3-3-1 models by weak charge data of Cesium and proton}, Nucl.\ Phys.\ B \textbf{943}, 114629 (2019), \href{https://doi.org/10.1016/j.nuclphysb.2019.114629}{doi:10.1016/j.nuclphysb.2019.114629}, \href{https://arxiv.org/abs/1812.08669}{arXiv:1812.08669}.
\bibitem{KEDR:2016hjj} V.~V.~Anashin \emph{et al.} (KEDR), \emph{Measurement of $R$ between 1.84 and 3.05 GeV at the KEDR detector}, Phys.\ Lett.\ B \textbf{770}, 174 (2017), \href{https://doi.org/10.1016/j.physletb.2017.04.073}{doi:10.1016/j.physletb.2017.04.073}, \href{https://arxiv.org/abs/1610.02827}{arXiv:1610.02827}.
\end{thebibliography}

\end{document}
